% =====================================================================
%  common.tex -- shared macros, theorem environments and book layout
%  Chapters may define their own local macros: every chapter is
%  wrapped in a TeX group (see \begin{notechapter}), so local
%  definitions never leak into other chapters.
% =====================================================================
\makeatletter

% ---------- helpers used by chapter-local definitions -----------------
% \forgetcs\foo     : forget \foo so that \newcommand\foo can (re)define it
% \forgetenv{foo}    : forget environment foo so it can be (re)defined
\newcommand*{\forgetcs}[1]{\let#1\relax}
\newcommand*{\forgetenv}[1]{%
  \expandafter\let\csname #1\endcsname\relax
  \expandafter\let\csname end#1\endcsname\relax}
\newcommand*{\guardnewlength}[1]{\ifdefined#1\else\newlength{#1}\fi}
\newcommand*{\guardnewcounter}[1]{\@ifundefined{c@#1}{\newcounter{#1}}{}}

% ---------- chapter wrapper ----------------------------------------------
\newenvironment{notechapter}{\begingroup}{\endgroup}

% Appendix inside a chapter (replaces \appendix of the original notes)
\newcommand{\chapterappendix}{%
  \par\setcounter{section}{0}%
  \renewcommand{\thesection}{\thechapter.\Alph{section}}%
  \renewcommand{\theHsection}{\thechapter.appendix.\Alph{section}}}

% Abstract (the book class has none)
\newenvironment{abstract}%
  {\begin{quote}\small\noindent\textbf{Abstract.}\ }%
  {\end{quote}\medskip}

% Per-chapter bibliography as an unnumbered section
\renewenvironment{thebibliography}[1]
  {\section*{References}%
   \list{\@biblabel{\@arabic\c@enumiv}}%
        {\settowidth\labelwidth{\@biblabel{#1}}%
         \leftmargin\labelwidth \advance\leftmargin\labelsep
         \@openbib@code \usecounter{enumiv}%
         \let\p@enumiv\@empty
         \renewcommand\theenumiv{\@arabic\c@enumiv}}%
   \sloppy\clubpenalty4000\@clubpenalty\clubpenalty
   \widowpenalty4000\sfcode`\.\@m}
  {\def\@noitemerr{\@latex@warning{Empty `thebibliography' environment}}%
   \endlist}

% Python blocks (originally pythontex `pyblock`) are typeset, not run
\lstnewenvironment{pyblock}{\lstset{language=Python}}{}

\makeatother

% ---------- default listings style ---------------------------------------
\definecolor{codegreen}{rgb}{0,0.6,0}
\definecolor{codegray}{rgb}{0.5,0.5,0.5}
\definecolor{codepurple}{rgb}{0.58,0,0.82}
\definecolor{backcolour}{rgb}{0.95,0.95,0.92}
\lstdefinestyle{mystyle}{
  backgroundcolor=\color{backcolour}, commentstyle=\color{codegreen},
  keywordstyle=\color{magenta}, numberstyle=\tiny\color{codegray},
  stringstyle=\color{codepurple}, basicstyle=\ttfamily\footnotesize,
  breakatwhitespace=false, breaklines=true, captionpos=b,
  keepspaces=true, numbers=left, numbersep=5pt, showspaces=false,
  showstringspaces=false, showtabs=false, tabsize=2}
\lstset{style=mystyle}

% ---------- theorem-like environments ------------------------------------
\theoremstyle{plain}
\newtheorem{theorem}{Theorem}[chapter]
\newtheorem{lemma}[theorem]{Lemma}
\newtheorem{proposition}[theorem]{Proposition}
\newtheorem{prop}[theorem]{Proposition}
\newtheorem{corollary}[theorem]{Corollary}
\newtheorem{conjecture}[theorem]{Conjecture}
\newtheorem{claim}[theorem]{Claim}
\newtheorem*{theorem*}{Theorem}
\theoremstyle{definition}
\newtheorem{definition}[theorem]{Definition}
\newtheorem{example}[theorem]{Example}
\newtheorem{assumption}[theorem]{Assumption}
\newtheorem{exercise}[theorem]{Exercise}
\newtheorem{notation}[theorem]{Notation}
\theoremstyle{remark}
\newtheorem{remark}[theorem]{Remark}
\newtheorem{note}[theorem]{Note}
\newtheorem{warning}[theorem]{Warning}
\newtheoremstyle{noparentheses}{\topsep}{\topsep}{\itshape}{0pt}{\bfseries}{.}
  {5pt plus 1pt minus 1pt}{\thmname{#1} \thmnumber{#2} \thmnote{#3}}

% cleveref names for theorem boxes created with \newtcbtheorem in some chapters
\crefname{tcb@cnt@definition}{Definition}{Definitions}
\crefname{tcb@cnt@theorem}{Theorem}{Theorems}
\crefname{tcb@cnt@proposition}{Proposition}{Propositions}
\crefname{tcb@cnt@lemma}{Lemma}{Lemmas}

% ---------- frequently used notation -------------------------------------
\newcommand{\R}{\mathbb{R}}
\newcommand{\N}{\mathbb{N}}
\newcommand{\Z}{\mathbb{Z}}
\newcommand{\C}{\mathbb{C}}
\newcommand{\E}{\mathbb{E}}
\newcommand{\Prob}{\mathbb{P}}
\newcommand{\bbone}{\mathbf{1}}
\DeclareMathOperator*{\argmax}{arg\,max}
\DeclareMathOperator*{\argmin}{arg\,min}
\DeclareMathOperator{\diag}{diag}
\DeclareMathOperator{\tr}{tr}
\DeclareMathOperator{\Tr}{Tr}
\DeclareMathOperator{\rank}{rank}
\DeclareMathOperator{\sign}{sign}
\DeclareMathOperator{\softmax}{softmax}
\newcommand{\abs}[1]{\left|#1\right|}
\newcommand{\norm}[1]{\left\|#1\right\|}
\newcommand{\inner}[2]{\left\langle #1, #2 \right\rangle}

% ---------- layout -------------------------------------------------------
\setcounter{secnumdepth}{2}
\setcounter{tocdepth}{1}
\numberwithin{equation}{chapter}
\numberwithin{figure}{chapter}
\numberwithin{table}{chapter}
\allowdisplaybreaks
\setlength{\parindent}{0pt}
\setlength{\parskip}{4pt plus 2pt}
\pagestyle{fancy}
\fancyhf{}
\fancyhead[LE,RO]{\thepage}
\fancyhead[RE]{\small\nouppercase{\leftmark}}
\fancyhead[LO]{\small\nouppercase{\rightmark}}
\renewcommand{\headrulewidth}{0.4pt}
\fancypagestyle{plain}{\fancyhf{}\fancyfoot[C]{\thepage}\renewcommand{\headrulewidth}{0pt}}

% ---------- support for chapters converted from Markdown (pandoc) ---------
\providecommand{\tightlist}{\setlength{\itemsep}{0pt}\setlength{\parskip}{0pt}}

% ---------- Unicode characters found in the sources (pdfLaTeX) ------------
\DeclareUnicodeCharacter{2192}{\ensuremath{\rightarrow}}
\DeclareUnicodeCharacter{03B8}{\ensuremath{\theta}}
\DeclareUnicodeCharacter{03C0}{\ensuremath{\pi}}
\DeclareUnicodeCharacter{2011}{-}
\DeclareUnicodeCharacter{2713}{\ding{51}}
\DeclareUnicodeCharacter{2717}{\ding{55}}
\DeclareUnicodeCharacter{2264}{\ensuremath{\leq}}
\DeclareUnicodeCharacter{2265}{\ensuremath{\geq}}
\DeclareUnicodeCharacter{2550}{=}
\DeclareUnicodeCharacter{2500}{-}
\lstset{inputencoding=utf8,extendedchars=true,literate=
  {π}{{$\pi$}}1 {θ}{{$\theta$}}1 {→}{{$\rightarrow$}}1 {×}{{$\times$}}1
  {✓}{{\ding{51}}}1 {≤}{{$\leq$}}1 {é}{{\'e}}1 {—}{{---}}1 {–}{{--}}1 {’}{{'}}1}

% ---------- chapter subtitle ------------------------------------------------
\newcommand{\chaptersubtitle}[1]{%
  \vspace*{-2.5em}{\large\itshape #1\par}\vspace{2em}}

% unique hyperref anchors for algorithmic line numbers
\makeatletter
\newcounter{algblockcount}
\AtBeginEnvironment{algorithmic}{\stepcounter{algblockcount}}
\providecommand{\theHALG@line}{}\renewcommand{\theHALG@line}{alg\arabic{algblockcount}.\arabic{ALG@line}}
\makeatother
